%% =========================================================================
%%  Аннотация (русская). Объём ≈1293 знака с пробелами при лимите 1500
%%  (Положение о ВКР МФТИ, Прил. 1 §3, не более 1500 знаков).
%%  Структура минимальна по Положению: цели/задачи, результаты, рекомендации
%%  (вводный абзац и блок ключевых слов опущены как не требуемые Положением).
%% =========================================================================

\chapter*{Аннотация}
\addcontentsline{toc}{chapter}{Аннотация}

\noindent
\textbf{Цель} --- построить и теоретически обосновать
квазиньютоновские обновления с памятью $p$ прошлых секущих
условий для задач математического программирования.

\noindent
\textbf{Задачи:} в общей мультисекантной параметризации
обновлений якобиана охарактеризовать ранг-1 случай,
сохраняющий $p$ прошлых секущих условий; построить алгоритм
SP-Broyden и доказать для него теоремы минимальности и
ограниченной деградации; построить симметричный аналог SS-U
и установить его условную Q-сверхлинейную сходимость;
провести численную верификацию.

\noindent
\textbf{Результаты.} В рамках мультисекантной параметризации
обновлений якобиана в ранг-1 классе построено обновление
SP-Broyden, точно сохраняющее $p{+}1$ секущих условий; для него
доказаны теоремы минимальности и ограниченной деградации
с проектором ранга $p{+}1$ (вместо ранга~1 у классического
Бройдена). Для симметричного случая построена конструкция
SS-U~--- ранг-1 коррекция в~$s_k^{\perp}$, сохраняющая текущее
секущее условие и симметрию, с условной Q-сверхлинейной
сходимостью по Dennis--More. Результаты подтверждены численно
на стандартных тестовых задачах.

\noindent
\textbf{Рекомендации.} На основании численных экспериментов:
SP-Broyden рекомендован вместо классического Бройдена в задачах,
где последний неустойчив или сходится медленно; SS-U~--- для
гладкой оптимизации в задачах, на которых SR1 теряет устойчивость.
